% Phần chứng minh chi tiết của Chương 2, được \input ở cuối 02-co-so-ly-thuyet.tex.
\section{Alpha-Beta: chứng minh đầy đủ và các hệ quả cho cài đặt}
\label{sec:lt-ab-day-du}
Mục~\ref{sec:lt-chung-minh} đã nêu ý chính vì sao cắt tỉa không đổi giá trị gốc. Cài đặt của đề tài còn dựa vào
những tính chất mạnh hơn: ý nghĩa của giá trị trả về khi nó nằm \emph{ngoài} cửa sổ (để ghi bảng chuyển vị), tính
đúng của cửa sổ lệch khi chỉ cần $k$ nước tốt nhất (để sinh dữ liệu), và quan hệ giữa độ sâu với số nút. Mục này
chứng minh từng tính chất, đối chiếu với đoạn mã tương ứng trong \repoFile{caro/agents/alphabeta.py}.

\subsection{Dạng Negamax}
Gọi $V_d(s)$ là giá trị Minimax độ sâu $d$ của thế $s$ \emph{theo góc nhìn bên đang đi}. Với $s$ kết thúc hoặc
$d=0$, $V_d(s)=E(s)$, trong đó $E$ là điểm thắng thua hoặc hàm lượng giá, cũng theo góc nhìn bên đang đi.

\begin{dinhly}[Tương đương Negamax]
\label{dl:cs-negamax}
Với trò chơi hai người tổng bằng không và quy ước góc nhìn ở trên,
\begin{equation}
  V_d(s)=\max_{a\in A(s)}\bigl(-V_{d-1}(T(s,a))\bigr),
  \label{eq:cs-negamax}
\end{equation}
và giá trị này trùng với Minimax viết theo góc nhìn cố định của một bên.
\end{dinhly}
\begin{proof}
Đặt $M_d(s)$ là giá trị Minimax theo góc nhìn bên X: khi X đi lấy lớn nhất, khi O đi lấy nhỏ nhất các
$M_{d-1}(T(s,a))$. Đặt $\sigma(s)=+1$ nếu X đi, $-1$ nếu O đi; ta chứng minh $M_d(s)=\sigma(s)V_d(s)$ bằng quy nạp theo $d$.
Với $d=0$ hoặc thế kết thúc, đẳng thức là định nghĩa của $E$. Giả sử đúng tới $d-1$. Vì lượt đi đổi sau mỗi nước,
$\sigma(T(s,a))=-\sigma(s)$. Nếu X đi:
\[
  M_d(s)=\max_a M_{d-1}(T(s,a))=\max_a \bigl(-V_{d-1}(T(s,a))\bigr)=V_d(s).
\]
Nếu O đi:
\begin{align*}
  M_d(s)&=\min_a M_{d-1}(T(s,a))=\min_a V_{d-1}(T(s,a))\\
        &=-\max_a\bigl(-V_{d-1}(T(s,a))\bigr)=-V_d(s).\qedhere
\end{align*}
\end{proof}
Tính tổng bằng không của $E$ là điều kiện cần: hàm lượng giá của engine là
$E(s)=S(\text{ta})-S(\text{địch})$ nên đổi bên đi thì đổi dấu. Biến thể A7 \texttt{no\_defense} chỉ lấy
$S(\text{ta})$, không còn đối xứng, nên với biến thể đó thuật toán không còn tính Minimax của một trò chơi tổng
bằng không nào; nó vẫn chạy được nhưng giá trị gốc chỉ mang nghĩa heuristic.

\subsection{Định lý fail-soft}
\label{sec:cs-fail-soft}
Cài đặt trả về \emph{điểm tốt nhất tìm được} kể cả khi điểm đó nằm ngoài cửa sổ $(\alpha,\beta)$ (kiểu
fail-soft). Định lý sau là cơ sở cho ba cờ \texttt{EXACT}, \texttt{LOWER}, \texttt{UPPER}.

\begin{dinhly}[Ý nghĩa của giá trị trả về]
\label{dl:cs-fail-soft}
Gọi $g=\mathrm{AB}(s,d,\alpha,\beta)$ với $\alpha<\beta$ (không dùng bảng chuyển vị). Khi đó:
\begin{enumerate}
  \item nếu $g\le\alpha$ thì $V_d(s)\le g$ (fail-low: $g$ là cận trên);
  \item nếu $\alpha<g<\beta$ thì $V_d(s)=g$ (giá trị chính xác);
  \item nếu $g\ge\beta$ thì $V_d(s)\ge g$ (fail-high: $g$ là cận dưới).
\end{enumerate}
\end{dinhly}
\begin{proof}
Quy nạp theo $d$. Ở lá, $g=V_d(s)$ nên cả ba mệnh đề đúng. Giả sử đúng ở độ sâu $d-1$. Tại nút $s$, các con
$c_1,c_2,\dots$ được duyệt theo thứ tự; trước khi duyệt $c_i$, cận dưới hiện tại là
$\alpha_i=\max(\alpha,v_1,\dots,v_{i-1})$, trong đó $v_j=-\mathrm{AB}(c_j,d-1,-\beta,-\alpha_j)$. Gọi
$u_i=V_{d-1}(c_i)$, nên $V_d(s)=\max_i(-u_i)$. Áp dụng giả thiết quy nạp cho con $c_i$ với cửa sổ
$(-\beta,-\alpha_i)$ và giá trị trả về $-v_i$:
\begin{itemize}
  \item nếu $v_i\ge\beta$ thì $-v_i\le-\beta$ (fail-low ở con), nên $u_i\le -v_i$, tức $-u_i\ge v_i$;
  \item nếu $v_i\le\alpha_i$ thì $-v_i\ge-\alpha_i$ (fail-high ở con), nên $u_i\ge -v_i$, tức $-u_i\le v_i$;
  \item nếu $\alpha_i<v_i<\beta$ thì $-u_i=v_i$.
\end{itemize}
\emph{Có cắt tỉa tại con $i$.} Khi đó $g=v_i\ge\beta$, và $V_d(s)\ge -u_i\ge v_i=g$: mệnh đề 3.

\emph{Không cắt tỉa.} Mọi $v_i<\beta$ và $g=\max_i v_i$. Nếu $g\le\alpha$ thì $\alpha_i=\alpha$ với mọi $i$,
nên $v_i\le\alpha_i$ và $-u_i\le v_i\le g$; suy ra $V_d(s)\le g$: mệnh đề 1. Nếu $\alpha<g<\beta$, gọi $j$ là
chỉ số đầu tiên có $v_j=g$. Vì mọi $v$ đứng trước nhỏ hơn $g$ và $\alpha<g$, ta có $\alpha_j<g=v_j<\beta$, do
đó $-u_j=g$. Với mọi $i$ khác, hoặc $v_i\le\alpha_i$ nên $-u_i\le v_i\le g$, hoặc $\alpha_i<v_i<\beta$ nên
$-u_i=v_i\le g$. Vậy $V_d(s)=g$: mệnh đề 2.
\end{proof}

Đoạn mã ghi bảng chuyển vị dịch nguyên ba mệnh đề này: cờ \texttt{EXACT} khi
$\alpha_0<g<\beta$ (với $\alpha_0$ là giá trị $\alpha$ lúc vào nút), \texttt{LOWER} khi $g\ge\beta$,
\texttt{UPPER} trong trường hợp còn lại. Khi đọc lại, một mục \texttt{LOWER} chỉ được dùng để nâng $\alpha$, một
mục \texttt{UPPER} chỉ để hạ $\beta$, và chỉ khi độ sâu đã lưu không nhỏ hơn độ sâu đang cần. Dùng mục sâu hơn
làm thay đổi đại lượng được tính ($V_{d'}$ thay cho $V_d$ với $d'>d$); đó là lý do định lý
\ref{dl:cs-fail-soft} được phát biểu cho bản không bảng chuyển vị, và phép thử A2 so giá trị gốc giữa hai cấu
hình bằng thực nghiệm.

\subsection{Cửa sổ lệch khi chỉ cần top-\texorpdfstring{$k$}{k}}
Hàm \texttt{scored\_moves} sinh nhãn policy cho cấp 5 bằng $k$ nước gốc tốt nhất. Chấm đầy đủ mọi nước thì tốn;
tuỳ chọn \texttt{prune\_root} tìm con thứ $i$ với cửa sổ $(-\infty,\,-\theta+1)$, trong đó $\theta$ là điểm hạng
$k$ trong số các nước đã chấm.
\begin{dinhly}
\label{dl:cs-topk}
Nếu mọi điểm là số nguyên thì tập $k$ nước điểm cao nhất (giải hoà theo thứ tự duyệt) và điểm của chúng không đổi
khi bật cắt tỉa gốc.
\end{dinhly}
\begin{proof}
Gọi $w_i=-V_{d-1}(c_i)$ là điểm thật của nước $i$ và $h$ là giá trị search con trả về, điểm ghi nhận là $-h$.
Cửa sổ của con là $(-\infty,-\theta+1)$. Nếu $h<-\theta+1$, tức $h\le-\theta$ do $h$ nguyên, thì theo
định lý~\ref{dl:cs-fail-soft} giá trị chính xác, nên $-h=w_i\ge\theta$. Nếu $h\ge-\theta+1$ thì
$V_{d-1}(c_i)\ge h$, nên $w_i\le -h\le\theta-1<\theta$: nước $i$ không vào top-$k$, và điểm ghi nhận $-h$ cũng nhỏ
hơn $\theta$ nên không chen vào danh sách. Vậy mọi nước thuộc top-$k$ được chấm chính xác, còn nước ngoài top-$k$
vẫn ở ngoài. Ngưỡng $\theta$ chỉ tăng nên lập luận đúng ở mọi bước.
\end{proof}
Một chi tiết của hàm này dẫn tới hệ quả lớn trong dữ liệu: danh sách được sắp xếp \emph{ổn định} theo điểm, và
các nước được duyệt theo chỉ số ô tăng dần. Khi mọi nước cùng điểm, $k$ nước ``tốt nhất'' chính là $k$ ứng viên
có chỉ số nhỏ nhất, tức nằm phía trên, bên trái. Mục~\ref{sec:cs-hoa-diem} đo ảnh hưởng của điều này lên tập dữ
liệu và lên mạng đã huấn luyện.

\subsection{Cây tối thiểu của Knuth--Moore}
\begin{dinhly}[Số lá khi thứ tự hoàn hảo]
\label{dl:cs-knuth-moore}
Trên cây đều hệ số rẽ nhánh $b$, độ sâu $d$, nếu nước tốt nhất luôn được duyệt đầu tiên thì Alpha-Beta lượng giá
đúng $b^{\lceil d/2\rceil}+b^{\lfloor d/2\rfloor}-1$ lá.
\end{dinhly}
\begin{proof}
Phân loại nút theo Knuth và Moore \cite{knuth1975alphabeta}. Nút loại 1 nằm trên biến chính: con đầu thuộc loại
1, $b-1$ con còn lại thuộc loại 2. Nút loại 2 (nút bị cắt): chỉ cần duyệt con đầu, con đó thuộc loại 3. Nút loại
3 (nút phải duyệt hết): cả $b$ con đều thuộc loại 2. Gọi $L_t(d)$ là số lá dưới một nút loại $t$ còn độ sâu $d$,
với $L_t(0)=1$. Ta có
\[
  L_1(d)=L_1(d-1)+(b-1)L_2(d-1),\qquad L_2(d)=L_3(d-1),\qquad L_3(d)=b\,L_2(d-1).
\]
Hai phương trình sau cho $L_2(d)=b\,L_2(d-2)$ với $L_2(0)=L_2(1)=1$, nên $L_2(d)=b^{\lfloor d/2\rfloor}$ và
$L_3(d)=b^{\lceil d/2\rceil}$. Thay vào phương trình đầu:
\[
  L_1(d)=1+(b-1)\sum_{k=0}^{d-1}b^{\lfloor k/2\rfloor}.
\]
Với $d=2m$, tổng bằng $2(b^m-1)/(b-1)$ nên $L_1=2b^m-1$. Với $d=2m+1$, thêm số hạng $b^m$ cho
$L_1=2b^m-1+(b-1)b^m=b^{m+1}+b^m-1$. Cả hai trường hợp trùng công thức đã nêu.
\end{proof}
Hệ số rẽ nhánh hiệu dụng của cây tối thiểu xấp xỉ $\sqrt b$. Với $b\approx 80$ (số ứng viên trung bình ở thế
trung cuộc), $\sqrt b\approx 9$. Bảng~\ref{tab:cs-tk-dosau} cho tỷ số số nút giữa hai độ sâu liên tiếp đo được
trên engine; chênh lệch so với 9 phản ánh ba nguồn: thứ tự chưa hoàn hảo, bảng chuyển vị gộp các chuyển vị (làm
giảm), và nước bắt buộc thu hẹp cây về một nhánh khi có đe doạ (cũng làm giảm).

\subsection{Chi phí của lặp sâu dần}
\begin{dinhly}
Nếu số nút ở độ sâu $d$ là $N(d)=c\,\rho^d$ với $\rho>1$ thì tổng số nút của lặp sâu dần tới độ sâu $D$ nhỏ hơn
$\dfrac{\rho}{\rho-1}\,N(D)$.
\end{dinhly}
\begin{proof}
$\sum_{d=1}^{D}c\rho^d = c\rho^D\sum_{k=0}^{D-1}\rho^{-k}<c\rho^D\dfrac{1}{1-\rho^{-1}}=\dfrac{\rho}{\rho-1}N(D)$.
\end{proof}
Với $\rho\approx9$, phần thêm tối đa khoảng $12{,}5\%$. Đổi lại, lượt tìm nông để lại thông tin sắp xếp (bảng
chuyển vị ở nút trong, lịch sử cắt tỉa) giúp lượt sâu hơn tiến gần cây tối thiểu, và ở chế độ thời gian luôn có
sẵn một nước hoàn chỉnh khi hết giờ. Riêng ở gốc, cài đặt hiện tại không bảo đảm duyệt nước tốt nhất của độ sâu
trước đầu tiên (Mục~\ref{sec:cs-goc}).

\subsection{Điểm thắng phụ thuộc ply và bảng chuyển vị}
\label{sec:cs-mate-tt}
Cài đặt cho điểm thế thua tại độ sâu $p$ tính từ gốc là $-(W-p)$ với $W=10^6$, nên giữa hai đường thắng, đường
ngắn hơn được chọn. Giá trị đó phụ thuộc $p$, trong khi khoá Zobrist chỉ phụ thuộc thế cờ. Một thế được lưu khi
gặp ở ply $p_1$ rồi đọc lại ở ply $p_2$ sẽ mang khoảng cách thắng lệch $p_1-p_2$. Vì mọi điểm thắng đều lớn hơn
mọi điểm lượng giá thường, \emph{dấu} và việc đây là thế thắng ép không đổi; chỉ thứ tự ưu tiên giữa các đường
thắng có thể sai. Cách sửa chuẩn là lưu $g+p$ (khoảng cách tính từ thế đang lưu) và trừ ply khi đọc. Mã hiện tại
chưa làm bước này; đây là một giới hạn đã biết của cấp 3, không ảnh hưởng kết quả thắng thua của một thế mà bộ
tìm kiếm đã chứng minh trong độ sâu của nó.

\subsection{Va chạm khoá Zobrist}
\label{sec:cs-zobrist}
Khoá của một thế là XOR các số 64 bit ngẫu nhiên ứng với từng quân và lượt đi \cite{zobrist1970}. Với các số độc
lập đều trên $\{0,1\}^{64}$, hai thế khác nhau cho khoá độc lập theo từng cặp, nên xác suất trùng là $2^{-64}$.
Với $m$ thế khác nhau được lưu, theo cận hợp
\begin{equation}
  P(\text{có ít nhất một va chạm})\le\binom{m}{2}2^{-64}<\frac{m^2}{2^{65}}.
  \label{eq:cs-va-cham}
\end{equation}
Với $m=10^6$ cận này khoảng $2{,}7\cdot10^{-8}$; với $m=10^9$ khoảng $0{,}027$. Bảng chuyển vị của đề tài giới
hạn $64$ MB, ước lượng $200$ byte mỗi mục, tức tối đa khoảng $3{,}4\cdot10^5$ mục, nên rủi ro va chạm không đáng
kể. Bảng không lưu lại thế cờ để kiểm tra, nên va chạm nếu xảy ra sẽ không bị phát hiện.

\section{Chuỗi ép bằng tứ: định nghĩa và tính đúng đắn}
\label{sec:cs-vcf-lt}
Chương~\ref{ch:thuc-nghiem} dùng một bộ giải chuỗi tứ liên tiếp (VCF, \emph{victory by continuous fours}) làm
trọng tài cho phân tích chiến thuật. Mục này phát biểu chính xác bài toán mà bộ giải giải và chứng minh lời giải
của nó là chắc chắn. Ý tưởng tìm kiếm trong không gian đe doạ được Allis và cộng sự dùng để giải Go-moku
\cite{allis1996gomoku,allis1994}; bộ giải ở đây chỉ xét lớp đe doạ mạnh nhất.

\subsection{Ký hiệu}
Một thế là cặp $(M,O)$ các tập ô của bên đang đi và đối thủ. $\mathcal W$ là tập các cửa sổ 5 ô liên tiếp theo
bốn hướng ($|\mathcal W|=252$ trên bàn 11$\times$11). Tập ô thắng ngay của bên $M$ trước $O$ là
\begin{equation}
  W(M,O)=\{\,e\notin M\cup O:\ \exists w\in\mathcal W,\ e\in w,\ w\cap O=\emptyset,\ |w\cap M|=4\,\}.
\end{equation}
Trong luật freestyle, đặt quân vào $e\in W(M,O)$ tạo đúng một dãy ít nhất 5 nên thắng ngay; ngược lại, một nước
thắng ngay luôn thuộc $W$.

\begin{dinhnghia}[Nước tứ và chuỗi VCF]
Nước $m\notin M\cup O$ là \emph{nước tứ} nếu $W(M\cup\{m\},O)\neq\emptyset$. Một \emph{chuỗi VCF độ dài $n$} từ
$(M,O)$ là:
\begin{itemize}
  \item $n=1$: một ô $e\in W(M,O)$;
  \item $n\ge2$: khi $W(M,O)=\emptyset$, một nước tứ $m$ sao cho $W(O,M\cup\{m\})=\emptyset$, và hoặc
  $|W(M\cup\{m\},O)|\ge2$ (khi đó $n=2$), hoặc $W(M\cup\{m\},O)=\{b\}$ và từ $(M\cup\{m\},\,O\cup\{b\})$ có chuỗi
  VCF độ dài $n-1$.
\end{itemize}
\end{dinhnghia}

\begin{dinhly}[Tính chắc chắn]
\label{dl:cs-vcf}
Nếu từ $(M,O)$, lượt của $M$, có chuỗi VCF độ dài $n$ thì bên $M$ có chiến lược thắng trong không quá $2n-1$ nước
đi (tính cả hai bên), bất kể đối thủ đáp thế nào.
\end{dinhly}
\begin{proof}
Quy nạp theo $n$. Với $n=1$, đặt quân vào $e\in W(M,O)$ thắng ngay. Với $n\ge2$, bên $M$ đi $m$. Vì
$W(O,M\cup\{m\})=\emptyset$, đối thủ không có nước thắng ngay. Xét đáp trả $r$ bất kỳ. Nếu
$r\notin W(M\cup\{m\},O)$ thì vẫn còn một ô $e\in W(M\cup\{m\},O)$, $e\ne r$; thêm quân đối thủ vào $r$ không
làm cửa sổ chứa $e$ có quân đối thủ trừ khi $r$ nằm trong cửa sổ đó, nhưng một cửa sổ có 4 quân $M$ và 1 ô trống
$e$ không còn ô nào khác để $r$ chiếm. Vậy $M$ đi $e$ và thắng ở nước thứ ba. Nếu $|W|\ge2$ và
$r\in W(M\cup\{m\},O)$, còn ít nhất một ô thắng khác và lập luận trên áp dụng: tổng 3 nước, đúng $2n-1$ với $n=2$.
Nếu $W=\{b\}$ và $r=b$, ta tới $(M\cup\{m\},O\cup\{b\})$ với lượt của $M$, theo giả thiết quy nạp thắng trong
$2(n-1)-1$ nước nữa, cộng hai nước vừa đi được $2n-1$.
\end{proof}

\paragraph{Điều kiện về phản đòn.} Nếu sau một nước chặn $b$, đối thủ tình cờ tạo được tứ của chính mình, thế mới
có $W(O',M')\neq\emptyset$. Định nghĩa buộc nước tứ tiếp theo của $M$ phải thoả $W(O',M'\cup\{m'\})=\emptyset$,
tức $m'$ vừa chặn tứ đối thủ vừa là tứ của mình; nếu đối thủ có từ hai ô thắng thì không nước nào thoả và nhánh bị
loại. Riêng khi $W(M',O')\neq\emptyset$ thì $M$ thắng ngay, không cần chặn.

\paragraph{Không đầy đủ.} Chiều ngược lại của định lý~\ref{dl:cs-vcf} sai: có thế thắng ép nhờ ba hở (VCT) hoặc
nhờ nước chuẩn bị yên tĩnh mà không có VCF. Vì vậy ``không có VCF'' không có nghĩa là thế an toàn; mọi số đo ở
Chương~\ref{ch:thuc-nghiem} dựa trên bộ giải phải đọc là \emph{cận dưới} của số đường thắng bị bỏ lỡ hay bị để hở.

\subsection{Đẳng thức tăng cường dùng khi sắp xếp nước}
Khi sắp xếp nước, cấp 3 cần $W(M\cup\{m\},O)$ cho mọi ứng viên $m$. Tính lại cho từng ứng viên tốn một lần quét
252 cửa sổ; đẳng thức sau cho phép tính một lần cho cả nút.
\begin{dinhly}
\label{dl:cs-gain}
Gọi $G(m)$ là tập các ô trống còn lại của những cửa sổ sống chứa đúng 3 quân $M$ và có $m$ là một trong hai ô
trống. Với mọi ô trống $m$:
\begin{equation}
  W(M\cup\{m\},O)=\bigl(W(M,O)\setminus\{m\}\bigr)\cup G(m).
  \label{eq:cs-gain}
\end{equation}
\end{dinhly}
\begin{proof}
($\supseteq$) Nếu $e\in W(M,O)$ và $e\ne m$, cửa sổ chứng tỏ điều đó có 4 quân và ô trống $e$; vì $m\ne e$ và $m$
trống nên $m$ nằm ngoài cửa sổ đó, và sau khi thêm $m$ cửa sổ vẫn nguyên. Nếu $e\in G(m)$, cửa sổ tương ứng có 3
quân và hai ô trống $m,e$; sau khi thêm $m$ nó có 4 quân và ô trống $e$.
($\subseteq$) Lấy $e\in W(M\cup\{m\},O)$ với cửa sổ $w$ có 4 quân của $M\cup\{m\}$. Rõ ràng $e\ne m$. Nếu
$m\notin w$ thì $w$ đã có 4 quân của $M$ nên $e\in W(M,O)\setminus\{m\}$. Nếu $m\in w$ thì $w$ có đúng 3 quân của
$M$ và hai ô trống $m,e$, nên $e\in G(m)$.
\end{proof}
Test \repoFile{tests/test_alphabeta.py} đối chiếu thứ tự nước thu được từ đẳng thức này với cách đi thử từng nước.

\section{Ước lượng Monte Carlo và giới hạn mẫu}
\label{sec:cs-mc}
\subsection{Rollout ước lượng cái gì}
Một rollout từ thế $s$ chơi đến hết ván theo chính sách mô phỏng $\pi_r$ và trả $Z\in\{-1,0,1\}$. Trung bình của
$n$ rollout độc lập là ước lượng không chệch của $\mu_{\pi_r}(s)=\mathbb E[Z\mid s,\pi_r]$, với phương sai
$\sigma^2/n$, $\sigma^2\le 1$. Đại lượng $\mu_{\pi_r}$ \emph{không phải} giá trị Minimax. Ví dụ: một thế đối thủ
có đúng một ô thắng ngay, bên đi có $k$ ứng viên. Giá trị Minimax không bị quyết định bởi đe doạ này (chỉ cần
chặn), nhưng rollout ngẫu nhiên chặn đúng với xác suất $1/k$; với $k=40$, thế bị đánh giá gần như thua. Rollout
\texttt{light} (thắng ngay, rồi chặn) sửa đúng lớp lỗi này nhưng vẫn bỏ qua đòn kép hai nước.

\subsection{Số mẫu cần để phân biệt hai nước}
\begin{dinhly}
\label{dl:cs-hoeffding}
Cho hai nước có kỳ vọng $\mu_1-\mu_2=\Delta>0$, phần thưởng trong $[-1,1]$, mỗi nước lấy $n$ mẫu độc lập. Nếu
$n\ge \dfrac{8}{\Delta^2}\ln\dfrac{4}{\delta}$ thì với xác suất ít nhất $1-\delta$, nước 1 có trung bình mẫu cao hơn.
\end{dinhly}
\begin{proof}
Bất đẳng thức Hoeffding \cite{hoeffding1963} cho biến trong khoảng độ dài 2:
$P(|\hat\mu-\mu|\ge\varepsilon)\le 2e^{-n\varepsilon^2/2}$. Lấy $\varepsilon=\Delta/2$:
$P\le 2e^{-n\Delta^2/8}$ cho mỗi nước; cận hợp cho hai nước là $4e^{-n\Delta^2/8}\le\delta$ khi $n$ thoả điều kiện.
Khi cả hai sai số nhỏ hơn $\Delta/2$, $\hat\mu_1>\mu_1-\Delta/2=\mu_2+\Delta/2>\hat\mu_2$.
\end{proof}
Với $\Delta=0{,}1$ và $\delta=0{,}05$ cần $n\ge 3\,506$ mẫu cho \emph{mỗi} nước. Chia đều 6\,400 mô phỏng cho 80
ứng viên chỉ được 80 mẫu mỗi nước, đủ để phân biệt chênh lệch cỡ $\Delta\approx0{,}66$. UCT không chia đều mà dồn
mẫu cho nước triển vọng, nhưng cận này giải thích vì sao ở ngân sách của đề tài, MCTS rollout chỉ phát hiện được
những khác biệt rất lớn; Bảng~\ref{tab:cs-ct-l4} đo trực tiếp điều đó.

\subsection{Chứng minh cận hồi quy của UCB1}
\begin{dinhly}[Auer, Cesa-Bianchi, Fischer \cite{auer2002ucb}]
Với $K$ máy có phần thưởng trong $[0,1]$, chiến lược chọn máy có $\bar X_a+\sqrt{2\ln t/n_a}$ lớn nhất cho
$\mathbb E[n_a(n)]\le \dfrac{8\ln n}{\Delta_a^2}+1+\dfrac{\pi^2}{3}$ với mọi máy dưới tối ưu $a$.
\end{dinhly}
\begin{proof}
Đặt $c_{t,s}=\sqrt{2\ln t/s}$ và $\ell=\lceil 8\ln n/\Delta_a^2\rceil$. Sau khi máy $a$ đã được kéo $\ell$ lần,
nó chỉ được chọn tại bước $t$ nếu tồn tại $s<t$ và $s_a\in[\ell,t)$ sao cho
$\bar X^*_s+c_{t,s}\le \bar X_{a,s_a}+c_{t,s_a}$. Bất đẳng thức này kéo theo ít nhất một trong ba biến cố:
\[
  \text{(i) } \bar X^*_s\le\mu^*-c_{t,s};\qquad \text{(ii) } \bar X_{a,s_a}\ge\mu_a+c_{t,s_a};\qquad
  \text{(iii) } \mu^*<\mu_a+2c_{t,s_a}.
\]
(Nếu cả ba sai thì $\bar X^*_s+c_{t,s}>\mu^*\ge\mu_a+2c_{t,s_a}>\bar X_{a,s_a}+c_{t,s_a}$.) Vì $s_a\ge\ell$ và
$t\le n$, $2c_{t,s_a}\le 2\sqrt{2\ln n/\ell}\le\Delta_a$, nên (iii) không xảy ra. Hoeffding cho mỗi biến cố (i),
(ii) xác suất không quá $e^{-4\ln t}=t^{-4}$. Cộng trên mọi $t,s,s_a$:
\[
  \mathbb E[n_a(n)]\le\ell+\sum_{t=1}^{\infty}\sum_{s=1}^{t}\sum_{s_a=1}^{t}2t^{-4}
  \le\ell+2\sum_{t=1}^{\infty}t^{-2}=\ell+\frac{\pi^2}{3},
\]
và $\ell\le 8\ln n/\Delta_a^2+1$.
\end{proof}
Nhân với $\Delta_a$ và cộng theo $a$ cho cận~\eqref{eq:ucb-regret}. Với phần thưởng trong $[-1,1]$, đổi biến
$r'=(r+1)/2$ cho thấy hằng số tương đương là $c_{uct}=2\sqrt2\approx2{,}83$. Đề tài dùng $c_{uct}=1{,}4$ (khai
thác nhiều hơn); Bảng~\ref{tab:cs-ct-cuct} đo ảnh hưởng của hằng số này trên bộ thế chiến thuật.

\section{Hành vi của quy tắc chọn PUCT}
\label{sec:cs-puct-lt}
Cấp 7 và 8 chọn con theo $Q(a)+c\,P(a)\sqrt{N}/(1+n_a)$ \cite{rosin2011puct,silver2017alphagozero}. Hai tính chất
sau đọc trực tiếp từ công thức và giải thích hành vi quan sát được.

\subsection{Phân bố lượt thăm tiệm cận}
Quy tắc tham lam trên điểm có xu hướng \emph{cân bằng điểm} giữa các nước đang cạnh tranh: nước nào điểm cao hơn
được thăm, làm số hạng thăm dò của nó giảm. Khi $N$ lớn và $Q$ đã ổn định, các nước được thăm đều có điểm xấp xỉ
một mức chung $\lambda$:
\begin{equation}
  Q(a)+c\,P(a)\frac{\sqrt N}{1+n_a}\approx\lambda
  \quad\Longrightarrow\quad
  1+n_a\approx \frac{c\sqrt N\,P(a)}{\lambda-Q(a)}\qquad(Q(a)<\lambda),
  \label{eq:cs-puct-can-bang}
\end{equation}
với $\lambda$ xác định bởi $\sum_a n_a=N$. Hai trường hợp giới hạn: khi mọi $Q$ bằng nhau, $n_a$ tỷ lệ với tiên
nghiệm $P(a)$; khi một nước có $Q$ vượt trội, $\lambda-Q$ của nó nhỏ và nó nhận phần lớn lượt thăm dù tiên nghiệm
không cao nhất. Grill và cộng sự chứng minh phân bố lượt thăm xấp xỉ nghiệm của bài toán tối ưu chính sách có
chính quy hoá $\max_y\,Q^\top y-\lambda_N\,\mathrm{KL}(P\,\Vert\,y)$ với $\lambda_N$ giảm theo cỡ $1/\sqrt N$
\cite{grill2020mcts}; nghiệm có đúng dạng~\eqref{eq:cs-puct-can-bang}.

Hệ quả quan trọng cho huấn luyện: nhãn policy $\pi_{\text{tìm kiếm}}\propto n_a$ ở ngân sách hữu hạn là một hỗn
hợp giữa tiên nghiệm và đánh giá tìm được. Ở ít mô phỏng, nhãn gần với chính mạng đang sinh ra nó, nên tín hiệu
cải thiện yếu.

\subsection{Giá trị của nước chưa thăm}
Mã đặt $Q=0$ cho con chưa được thăm (``coi như hoà''). Khi thế đang tốt (các con đã thăm có $Q>0$), con chưa thăm
kém hấp dẫn, tìm kiếm tập trung; khi thế đang xấu ($Q<0$), mọi con chưa thăm trông tốt hơn các con đã thăm, tìm
kiếm tự lan rộng. Đây là một lựa chọn thiết kế (first-play urgency); các lựa chọn khác như lấy $Q$ của cha chưa
được thử trong đề tài.

\subsection{Virtual loss}
Để gom $B$ lá cho một lần suy luận, mỗi lượt đi xuống cộng một ``trận thua tạm'' vào các nút trên đường: số thăm
hiệu dụng $n+v$ và tổng giá trị của con tăng $v$ (con thắng, tức cha thua). Điểm của đường vừa đi giảm ở cả hai
số hạng, nên lượt kế tiếp rẽ sang nhánh khác. Khi $B=1$, virtual loss không còn tồn tại lúc chọn nên PUCT trùng
bản tuần tự. Khi $B$ lớn, các lựa chọn trong lô dựa trên thông tin cũ và bị đẩy ra khỏi nhánh tốt; đây là đánh
đổi giữa thông lượng suy luận và chất lượng tìm kiếm, đo ở Bảng~\ref{tab:cs-ct-batch} và
\ref{tab:cs-tk-backend-puct}.

\section{Mạng nơ-ron: các tính chất dùng trong phân tích}
\label{sec:cs-mang-lt}
\subsection{Trường tiếp nhận}
Mỗi lớp tích chập $3\times3$ đệm 1 mở rộng trường tiếp nhận thêm 2 ô. CaroResNet có một lớp vào và sáu khối, mỗi
khối hai lớp, tổng 13 lớp, nên trường tiếp nhận là $1+2\cdot13=27$ ô. Để một ô ở góc ``nhìn thấy'' góc đối diện
trên bàn 11$\times$11 cần $21$ ô, tức 10 lớp. Vậy từ khối thứ năm trở đi, mỗi vị trí đầu ra phụ thuộc toàn bàn;
các đầu ra không bị giới hạn bởi tính cục bộ.

\subsection{Đẳng biến tịnh tiến và đối xứng}
Tích chập đẳng biến với phép tịnh tiến (ở xa biên), nhưng không đẳng biến với phép quay hay lật: không có lý do
cấu trúc nào để $f(g\cdot x)=g\cdot f(x)$ với $g\in D_4$. Đầu policy còn có lớp kết nối đầy đủ
$242\to121$ gán trọng số riêng cho từng ô, nên cũng mất đẳng biến tịnh tiến. Vì vậy đối xứng của luật chơi phải
được \emph{học}, qua tăng cường dữ liệu, hoặc \emph{áp đặt}, qua lấy trung bình trên nhóm.

\begin{dinhly}
\label{dl:cs-d4}
Cho $f$ trả phân bố trên các ô. Đặt $\bar f(x)=\frac18\sum_{g\in D_4}g^{-1}\cdot f(g\cdot x)$. Khi đó
$\bar f(h\cdot x)=h\cdot\bar f(x)$ với mọi $h\in D_4$; và với mọi nhãn $\pi$, entropy chéo của $\bar f(x)$ không
lớn hơn trung bình entropy chéo của tám phân bố $g^{-1}\cdot f(g\cdot x)$.
\end{dinhly}
\begin{proof}
Đổi biến $k=gh$ (khi $g$ chạy khắp nhóm, $k$ cũng vậy), $g=kh^{-1}$, $g^{-1}=hk^{-1}$:
$\bar f(h\cdot x)=\frac18\sum_g g^{-1}f(gh\cdot x)=\frac18\sum_k h k^{-1}f(k\cdot x)=h\cdot\bar f(x)$.
Phần hai: $-\log$ lồi nên theo Jensen
$-\sum_a\pi_a\log\frac18\sum_g p_g(a)\le \frac18\sum_g\bigl(-\sum_a\pi_a\log p_g(a)\bigr)$.
\end{proof}
Định lý chỉ so với \emph{trung bình} của tám hướng, không bảo đảm tốt hơn hướng gốc; việc đó là kết quả thực
nghiệm ở Bảng~\ref{tab:cs-mang-doixung}.

\subsection{Gradient của hàm mất mát}
Gọi $z\in\mathbb R^{121}$ là logit, $p=\mathrm{softmax}(z)$, $\pi$ là nhãn. Với $L_p=-\sum_a\pi_a\log p_a$ và
$\log p_a=z_a-\log\sum_b e^{z_b}$:
\begin{equation}
  \frac{\partial L_p}{\partial z_b}=-\pi_b+\Bigl(\sum_a\pi_a\Bigr)p_b=p_b-\pi_b.
\end{equation}
Gradient bị chặn và bằng 0 đúng khi $p=\pi$. Với đầu value $v=\tanh u$ và $L_v=(v-z)^2$:
\begin{equation}
  \frac{\partial L_v}{\partial u}=2(v-z)(1-v^2).
\end{equation}
Khi $v$ bão hoà sai phía (gần $-1$ mà $z=1$), thừa số $1-v^2$ gần 0 làm gradient tắt. Đây là một lý do lỗi quy ước
dấu ở SL v2 (Chương~\ref{ch:phuong-phap}) khó tự sửa trong học tăng cường.

\subsection{Chuẩn hoá theo lô và kết nối tắt}
Batch normalization biến đổi mỗi kênh $x\mapsto\gamma\,(x-\mu_B)/\sqrt{\sigma_B^2+\epsilon}+\beta$ với thống kê
của lô khi huấn luyện, và thống kê chạy tích luỹ khi suy luận. Vì vậy cùng một thế cho đầu ra khác nhau ở hai chế
độ; mọi tác tử gọi \texttt{model.eval()} trước khi suy luận. Kết nối tắt $y=\mathrm{ReLU}(x+F(x))$ \cite{he2016resnet}
cho $\partial(x+F(x))/\partial x=I+\partial F/\partial x$, giữ một đường truyền gradient không suy giảm qua sáu khối.

\subsection{Calibration và phân rã sai số value}
\begin{dinhly}
\label{dl:cs-mse}
Với $v=v(S)$ dự đoán và $z$ kết quả,
$\mathbb E[(v-z)^2]=\mathbb E\bigl[(v-\mathbb E[z\mid S])^2\bigr]+\mathbb E\bigl[\mathrm{Var}(z\mid S)\bigr]$.
\end{dinhly}
\begin{proof}
Viết $v-z=(v-m)+(m-z)$ với $m=\mathbb E[z\mid S]$ rồi bình phương. Hạng chéo
$\mathbb E[(v-m)(m-z)]=\mathbb E\bigl[(v-m)\,\mathbb E[m-z\mid S]\bigr]=0$ vì $v,m$ là hàm của $S$.
\end{proof}
Số hạng thứ hai là nhiễu không giảm được: cùng một thế có thể dẫn tới kết quả khác nhau do người chơi. Vì vậy MSE
không thể về 0, và hằng số 0 cho MSE bằng $\mathbb E[z^2]$ (tỷ lệ ván có thắng thua). Mạng được gọi là calibrated
nếu $\mathbb E[z\mid v]=v$; Hình~\ref{fig:cs-mang-calib} vẽ đại lượng vế trái theo $v$.

\section{Định lý policy gradient cho ván hai người}
\label{sec:cs-pg-lt}
\subsection{Dạng quỹ đạo}
Một ván tự đấu cho quỹ đạo $\tau=(s_0,a_0,\dots,s_T)$ với xác suất
$P_\theta(\tau)=\prod_t\pi_\theta(a_t\mid s_t)$ (chuyển trạng thái tất định; khai cuộc ngẫu nhiên không phụ thuộc
$\theta$). Gọi $z_t\in\{-1,0,1\}$ là kết quả theo góc nhìn bên đi ở bước $t$. Với mục tiêu của một bên,
\begin{equation}
  \nabla_\theta\,\mathbb E_\tau[R(\tau)]
  =\mathbb E_\tau\Bigl[R(\tau)\sum_t\nabla_\theta\log\pi_\theta(a_t\mid s_t)\Bigr],
\end{equation}
vì $\nabla P_\theta=P_\theta\nabla\log P_\theta$. Khi học cả hai bên bằng cùng tham số, mỗi nước dùng $z_t$ của
chính bên đi nước đó; đó là dòng tính \texttt{zs} trong \repoFile{caro/training/train_pg.py}. Gradient của nước
$t$ chỉ cần phần thưởng về sau, và ở đây phần thưởng duy nhất ở cuối ván nên $z_t$ dùng trực tiếp.

\subsection{Baseline tối ưu và điều kiện giảm phương sai}
Đặt $g=\nabla_\theta\log\pi_\theta(a\mid s)$. Ước lượng $(z-b)g$ không chệch với mọi $b=b(s)$. Phương sai tổng
(vết hiệp phương sai) là $\mathbb E[(z-b)^2\Vert g\Vert^2]-\Vert\mathbb E[zg]\Vert^2$; lấy đạo hàm theo $b$ tại
mỗi $s$ cho baseline tối ưu
\begin{equation}
  b^*(s)=\frac{\mathbb E\bigl[z\,\Vert g\Vert^2\mid s\bigr]}{\mathbb E\bigl[\Vert g\Vert^2\mid s\bigr]}.
\end{equation}
Value head xấp xỉ $\mathbb E[z\mid s]$, trùng $b^*$ khi $\Vert g\Vert^2$ ít tương quan với $z$. Nếu coi
$\Vert g\Vert^2$ gần hằng, baseline $v$ giảm phương sai khi và chỉ khi $\mathbb E[(z-v)^2]<\mathbb E[z^2]$. Với lỗi
quy ước dấu đã phát hiện, MSE của value theo nhãn bên đi là $1{,}38$ trong khi hằng số 0 đạt $0{,}91$
(\repoFile{docs/GHI_CHU_value_head.md}): baseline khi đó \emph{làm tăng} phương sai.

\subsection{Vai trò của entropy và nhiệt độ}
Thưởng entropy $\beta H(\pi)$ có gradient theo logit $\beta\,\partial H/\partial z_b=-\beta\,p_b(\log p_b+H)$, đẩy
xác suất cao xuống và xác suất thấp lên. Khi lấy mẫu với nhiệt độ $\tau$, $\pi=\mathrm{softmax}(z/\tau)$ và
$\partial\log\pi_a/\partial z_b=(\mathbb 1[a{=}b]-\pi_b)/\tau$; đề tài dùng $\tau=1$ và $\beta=0{,}01$ giảm theo
hệ số $0{,}999$ mỗi vòng, tức sau 390 vòng còn $\beta\approx0{,}0068$.

\subsection{Cổng chọn mạng như một phép thử}
Vòng AlphaZero chỉ nhận mạng mới khi điểm của nó trước mạng đương nhiệm đạt $\theta=0{,}55$ trên $n=100$ ván
\cite{silver2017alphagozero}. Nếu các ván độc lập và điểm thật là $p$ (bỏ qua hoà),
$P(\text{nhận})\approx\Phi\bigl((p-\theta)/\sqrt{p(1-p)/n}\bigr)$, cho trong Bảng~\ref{tab:cs-gating}.
\begin{table}[H]
  \centering\small
  \caption{Xác suất xấp xỉ một mạng ứng viên qua cổng $0{,}55$ với 100 ván độc lập}
  \label{tab:cs-gating}
  \begin{tabular}{lrrrrr}
    \toprule
    Điểm thật $p$ & 0{,}45 & 0{,}50 & 0{,}55 & 0{,}60 & 0{,}65\\
    \midrule
    Xác suất nhận (\%) & 2{,}2 & 15{,}9 & 50{,}0 & 84{,}6 & 98{,}2\\
    \bottomrule
  \end{tabular}
\end{table}
Một mạng không mạnh hơn vẫn qua cổng khoảng một trong sáu lần. Các ván của cổng chơi tất định ($\tau=0$) trên
khai cuộc đổi màu nên không độc lập hoàn toàn; bảng này chỉ cho bậc độ lớn.

\section{Chi tiết các phương pháp thống kê}
\label{sec:cs-tk-lt}
\subsection{Kiểm định dấu theo khai cuộc và kiểm định McNemar}
Hai ván đổi màu trên một khai cuộc cho điểm gộp thuộc $\{0;\,0{,}25;\,0{,}5;\,0{,}75;\,1\}$. Dưới giả thuyết hai tác
tử ngang sức, khai cuộc ``nghiêng'' về mỗi bên với xác suất bằng nhau, nên số khai cuộc nghiêng về A trong số
$n$ khai cuộc không hoà điểm theo phân bố $\mathrm{Bin}(n,\tfrac12)$. Giá trị $p$ hai phía chính xác là
$\min\bigl(1,\,2\sum_{k\ge\max(w,l)}\binom nk 2^{-n}\bigr)$. Kiểm định này không cần giả định phân bố điểm và dùng
đúng đơn vị độc lập (khai cuộc). Kiểm định McNemar \cite{mcnemar1947} áp dụng cùng công thức cho hai tác tử giải
cùng một tập thế: chỉ những thế một bên đúng, một bên sai mang thông tin.

\subsection{Hiệu chỉnh Holm}
Sắp $m$ giá trị $p_{(1)}\le\dots\le p_{(m)}$, đặt
$\tilde p_{(i)}=\max_{j\le i}\min\bigl(1,(m-j+1)p_{(j)}\bigr)$ và bác bỏ khi $\tilde p_{(i)}<\alpha$
\cite{holm1979}.
\begin{dinhly}
Thủ tục Holm kiểm soát xác suất bác bỏ nhầm ít nhất một giả thuyết đúng (FWER) ở mức $\alpha$, không cần giả định
độc lập.
\end{dinhly}
\begin{proof}
Gọi $I$ là tập giả thuyết đúng, $m_0=|I|$, và $(j)$ là vị trí đầu tiên trong dãy đã sắp thuộc $I$. Có ít nhất
$j-1$ giả thuyết sai đứng trước nên $j\le m-m_0+1$, tức $m-j+1\ge m_0$. Muốn bác bỏ nhầm phải bác bỏ tới vị trí
$j$, tức $(m-j+1)p_{(j)}<\alpha$, kéo theo $p_{(j)}<\alpha/m_0$. Vậy biến cố bác bỏ nhầm nằm trong
$\bigcup_{i\in I}\{p_i<\alpha/m_0\}$, có xác suất không quá $m_0\cdot\alpha/m_0=\alpha$.
\end{proof}

\subsection{Bootstrap phân vị}
Với đơn vị quan sát $x_1,\dots,x_m$ (ở đây là khai cuộc) và thống kê $T$, bootstrap \cite{efron1979} lấy $B$ mẫu
có hoàn lại cỡ $m$, tính $T^*_b$ và dùng phân vị $\alpha/2$, $1-\alpha/2$ của $\{T^*_b\}$. Phương pháp thay phân
bố chưa biết của dữ liệu bằng phân bố thực nghiệm. Nó không tạo thêm thông tin: với $m=50$ khai cuộc, khoảng vẫn
rộng, và khi mọi khai cuộc cho cùng kết quả (cặp toàn thắng), khoảng bootstrap suy biến thành một điểm. Trong
trường hợp đó khoảng Wilson trên số ván là cách trình bày thận trọng hơn.

\subsection{Cỡ mẫu}
Để phát hiện $p_1\neq\tfrac12$ bằng kiểm định hai phía mức $\alpha$ với lực $1-\beta$, xấp xỉ chuẩn cho
\begin{equation}
  n\ge\left(\frac{z_{1-\alpha/2}\cdot\frac12+z_{1-\beta}\sqrt{p_1(1-p_1)}}{p_1-\frac12}\right)^2,
\end{equation}
vì thống kê $\hat p$ có độ lệch chuẩn $\tfrac12/\sqrt n$ dưới giả thuyết không và $\sqrt{p_1(1-p_1)/n}$ dưới giả
thuyết đối; điều kiện lực là khoảng cách $p_1-\tfrac12$ bằng tổng hai biên. Giá trị tính sẵn ở
Bảng~\ref{tab:cs-tk-luc}.

\subsection{Bradley-Terry: tồn tại, duy nhất và cách giải}
\begin{dinhly}
\label{dl:cs-bt}
Đặt $\theta_i$ là điểm trên thang log-odds, $\ell(\theta)=\sum_{i\ne j}y_{ij}\log\sigma(\theta_i-\theta_j)$. Nếu mỗi
cặp đã gặp nhau có $y_{ij}>0$ và $y_{ji}>0$, và đồ thị các cặp đã gặp liên thông, thì $\ell$ có điểm cực đại duy
nhất trên mặt phẳng $\sum_i\theta_i=0$.
\end{dinhly}
\begin{proof}
Hessian của $\ell$ là $-L$, trong đó $L$ là Laplacian của đồ thị với trọng số cạnh
$w_{ij}=n_{ij}p_{ij}(1-p_{ij})>0$, $n_{ij}=y_{ij}+y_{ji}$. Với mọi $x$, $x^\top Lx=\sum_{i<j}w_{ij}(x_i-x_j)^2\ge0$,
bằng 0 chỉ khi $x$ hằng trên mỗi thành phần liên thông. Đồ thị liên thông nên trên mặt phẳng tổng bằng 0, $\ell$
lõm chặt: cực đại nếu có là duy nhất. Tồn tại: nếu $\Vert\theta\Vert\to\infty$ trên mặt phẳng đó thì
$\max_i\theta_i-\min_i\theta_i\to\infty$, nên theo tính liên thông có một cạnh với $\theta_i-\theta_j\to-\infty$ mà
$y_{ij}>0$; số hạng đó tiến tới $-\infty$ và các số hạng khác không dương, nên $\ell\to-\infty$. Hàm liên tục tiến
tới $-\infty$ ở vô cực đạt cực đại.
\end{proof}
Ván hoà ảo bảo đảm điều kiện $y_{ij},y_{ji}>0$ ngay cả khi một cặp toàn thắng; Ford chỉ ra điều kiện liên thông
mạnh tương tự cho dữ liệu gốc \cite{ford1957}, và Hunter bàn các thuật toán MM thay thế \cite{hunter2004}. Mã giải
bằng Newton: bước $\Delta=(L+\tfrac1n\mathbf 1\mathbf 1^\top)^{-1}\nabla\ell$; số hạng hạng một làm ma trận khả
nghịch mà không đổi thành phần vuông góc với $\mathbf 1$, rồi chia đôi bước cho tới khi $\ell$ không giảm.
